\documentclass[12pt]{article}
\usepackage{amsmath,amssymb}

\begin{document}
\section*{{2022 AIME I Problems/Problem 3}}
Source: AoPS Wiki

\subsection*{Solution}
We have the following diagram: [asy] /* Made by MRENTHUSIASM , modified by Cytronical */ size(300); pair A, B, C, D, A1, B1, C1, D1, P, Q, X, Y, Z, W; A = (-250,6*sqrt(731)); B = (250,6*sqrt(731)); C = (325,-6*sqrt(731)); D = (-325,-6*sqrt(731)); A1 = bisectorpoint(B,A,D); B1 = bisectorpoint(A,B,C); C1 = bisectorpoint(B,C,D); D1 = bisectorpoint(A,D,C); P = intersectionpoint(A--300*(A1-A)+A,D--300*(D1-D)+D); Q = intersectionpoint(B--300*(B1-B)+B,C--300*(C1-C)+C); X = intersectionpoint(B--5*(Q-B)+B,C--D); Y = (0,6*sqrt(731)); Z = intersectionpoint(A--4*(P-A)+A,B--4*(Q-B)+B); W = intersectionpoint(A--5*(P-A)+A,C--D); draw(anglemark(P,A,B,1000),red); draw(anglemark(D,A,P,1000),red); draw(anglemark(A,B,Q,1000),red); draw(anglemark(Q,B,C,1000),red); draw(anglemark(P,D,A,1000),red); draw(anglemark(C,D,P,1000),red); draw(anglemark(Q,C,D,1000),red); draw(anglemark(B,C,Q,1000),red); add(pathticks(anglemark(P,A,B,1000), n = 1, r = 0.15, s = 750, red)); add(pathticks(anglemark(D,A,P,1000), n = 1, r = 0.15, s = 750, red)); add(pathticks(anglemark(A,B,Q,1000), n = 1, r = 0.15, s = 750, red)); add(pathticks(anglemark(Q,B,C,1000), n = 1, r = 0.15, s = 750, red)); add(pathticks(anglemark(P,D,A,1000), n = 2, r = 0.12, spacing = 150, s = 750, red)); add(pathticks(anglemark(C,D,P,1000), n = 2, r = 0.12, spacing = 150, s = 750, red)); add(pathticks(anglemark(Q,C,D,1000), n = 2, r = 0.12, spacing = 150, s = 750, red)); add(pathticks(anglemark(B,C,Q,1000), n = 2, r = 0.12, spacing = 150, s = 750, red)); dot("$A$",A,1.5*dir(A),linewidth(4)); dot("$B$",B,1.5*dir(B),linewidth(4)); dot("$C$",C,1.5*dir(C),linewidth(4)); dot("$D$",D,1.5*dir(D),linewidth(4)); dot("$P$",P,1.5*(-1,0),linewidth(4)); dot("$Q$",Q,1.5*E,linewidth(4)); dot("$X$",X,1.5*dir(-105),linewidth(4)); dot("$Y$",Y,1.5*N,linewidth(4)); dot("$Z$",Z,4.5*dir(75),linewidth(4)); dot("$W$",W,1.5*dir(-75),linewidth(4)); draw(A--B--C--D--cycle^^A--P--D^^B--Q--C^^P--Q); draw(P--W^^Q--X^^Y--Z,dashed); [/asy] Let $X$ and $W$ be the points where $AP$ and $BQ$ extend to meet $CD$ , and $YZ$ be the height of $\triangle AZB$ . As proven in Solution 2, triangles $APD$ and $DPW$ are congruent right triangles. Therefore, $AD = DW = 333$ . We can apply this logic to triangles $BCQ$ and $XCQ$ as well, giving us $BC = CX = 333$ . Since $CD = 650$ , $XW = DW + CX - CD = 16$ . Additionally, we can see that $\triangle XZW$ is similar to $\triangle PQZ$ and $\triangle AZB$ . We know that $\frac{XW}{AB} = \frac{16}{500}$ . So, we can say that the height of the triangle $AZB$ is $500u$ while the height of the triangle $XZW$ is $16u$ . After that, we can figure out the distance from $Y$ to $PQ: \frac{500+16}{2} = 258u$ and the height of triangle $PZQ: 500-258 = 242u$ . Finally, since the ratio between the height of $PZQ$ to the height of $AZB$ is $242:500$ and $AB$ is $500$ , $PQ = \boxed{242}.$ ~Cytronical

\end{document}
